\SourceLeaf{45}
\[ \cos a \cos b = \tfrac12\cos(a-b)+\tfrac12\cos(a+b). \]

\par\noindent\hspace*{\parindent}Sumatur iam $\cos\varphi$, quae, in $\cos\varphi$ ducta, dabit
$\cos\varphi . \cos\varphi$, quod productum tum, si revera multiplicemus,
per $\cos\varphi^2$, tum si per formulam modo expositam in summam
solvamus, per $\tfrac12+\tfrac12\cos 2\varphi$ exhibebitur. Est igitur
$\cos\varphi^2 = \tfrac12+\tfrac12\cos 2\varphi$. Simili ratione ad
potestates altiores provecti, sensim formulas eruimus sequentes

\[ \cos\varphi^3 = \tfrac34\cos\varphi + \tfrac14\cos 3\varphi ; \]
\[ \cos\varphi^4 = \tfrac38 + \tfrac48\cos 2\varphi + \tfrac18\cos 4\varphi ; \]
\[ \cos\varphi^5 = \tfrac{10}{16}\cos\varphi + \tfrac{5}{16}\cos 3\varphi + \tfrac{1}{16}\cos 5\varphi ; \]
\[ \cos\varphi^6 = \tfrac{10}{32} + \tfrac{15}{32}\cos 2\varphi + \tfrac{6}{32}\cos 4\varphi + \tfrac{1}{32}\cos 6\varphi ; \]
\[ \cos\varphi^7 = \tfrac{35}{64}\cos\varphi + \tfrac{21}{64}\cos 3\varphi + \tfrac{7}{64}\cos 5\varphi + \tfrac{1}{64}\cos 7\varphi ; \]
\[ \cos\varphi^8 = \tfrac{35}{128} + \tfrac{56}{128}\cos 2\varphi + \tfrac{28}{128}\cos 4\varphi + \tfrac{8}{128}\cos 6\varphi + \tfrac{1}{128}\cos 8\varphi . \]

\par\noindent Facile ex his calculis lex generalis eruitur. In potestatibus $\cos\varphi$
paribus non nisi parium ipsius $\varphi$ multiplorum, quorum maximum dato
exponente definitur, cosinus adparent, in imparibus eadem lege multiplorum
imparium cosinus. Sed pares non eandem omnino quam impares legem sequuntur.

$\alpha$. Terminus initialis in paribus, in $\cos 0\varphi = 1$ ductus,
coefficiente gaudet, cuius denominator, ut ceterorum omnium, est ipsius 2
potestas, cuius exponens unitate exceditur ab exponente propositae
potestatis. Numerator est fere coefficiens binomialis ad potestatem
propositam pertinens, et quidem, si illos suo ordine evolutos concipias,
mediorum inter illos dimidius. Cetera membra, multiplum ipsius $\varphi$,
indicis membri duplum, in coefficiente adiecto eundem quem initiale
denominatorem, pro numeratore autem coefficientem binomialem eiusdem
potestatis, illo quem initiale gerit tot gradibus posteriorem, quot index
membri indicat, continent. Quae lex ita generatim exponi potest: fore
$\cos\varphi^{2m}$ seriem, cuius terminus initialis\HistoricalNote{*)}{*)}{Expressio legis huius \textsc{Euler}iana, ad coefficientes
binomiales non revocata, cum proposita si ipsa congruit, sed non aeque
concinna videtur.

Ceterum coefficientes binomiales a me semper more \textsc{Hindenburg}iano
notari, et serierum terminos, duce Ill.\ \textsc{Kaestnero}, semper in
initialem et sequentes distingui, vix est quod moneam.}

\[ \frac{1}{2^{2m-1}} . \frac{{}^{2m}\mathfrak{B}^{(m)}}{2}, \]
